\documentclass{article}
\usepackage[utf8]{inputenc}
\usepackage{amsmath,amssymb}
\usepackage[most]{tcolorbox}
\usepackage{geometry}
\geometry{margin=2.5cm}

\newcommand{\step}[1]{\par\noindent\hangindent=3.5em\hangafter=1%
  \makebox[3.5em][l]{\textbf{#1.}}}
\newcommand{\steptag}[1]{\hfill{\small\textsf{[#1]}}}

\newtcolorbox{proofbox}[1]{
  colback=gray!5, colframe=gray!60,
  fonttitle=\bfseries, title={#1},
  breakable, enhanced,
  left=4pt, right=4pt, top=4pt, bottom=4pt,
  before skip=10pt, after skip=10pt
}

\begin{document}

\begin{proofbox}{Parameterized polar-path transport: there exist C\_path\textasciicircum{}0,...}
\step{1} Parameterized polar-path transport: there exist C\_path\textasciicircum{}0, C\_pol\textasciicircum{}0 >= 1 and kappa\_pol\textasciicircum{}0 in (0, 1/2] such that, for every def-stage1-polar-witness-data tuple W with C\_path >= C\_path\textasciicircum{}0, C\_pol >= C\_pol\textasciicircum{}0, and 0 < kappa\_pol <= kappa\_pol\textasciicircum{}0, for every finite-dimensional exact-unit epsilon\_r-C*-algebra, every delta > 0 with C\_pol*(epsilon\_r + delta) <= kappa\_pol, every U\_0, U\_1 in calU, and every q in [0, 1] satisfying ||U\_1 - U\_0|| <= q, C\_path*q <= 1/4, and C\_path*(q + epsilon\_r*q + q\textasciicircum{}2) < delta - C\_pol*(epsilon\_r*delta + delta\textasciicircum{}2), every L\_\{Z\_t\} is invertible and every Z\_t = (1-t)*U\_0 + t*U\_1 lies in calUbar\_\{C\_path*(q + epsilon\_r*q + q\textasciicircum{}2)\} for t in [0, 1], and, writing u\_delta for the unique first component of the inverse of Pi\_delta: calU x B\textasciicircum{}\{calH\}\_delta(J) -> S\_delta := Pi\_delta(calU x B\textasciicircum{}\{calH\}\_delta(J)), Pi\_delta(U, H) = U bold-dot H, the map H(t, U\_0, U\_1) = u\_delta(Z\_t) is jointly continuous in (t, U\_0, U\_1), joins U\_0 to U\_1, and satisfies H(t, cU\_0, cU\_1) = c*H(t, U\_0, U\_1) for every c in U(1). \steptag{validated} \\[6pt]
\step{1.1} Invoke lem-stage1-polar-path-admissibility and choose its universal witnesses A,B >= 1 and k in (0,1/2]. Invoke lem-stage1-polar-retraction-transport and choose its universal witnesses P >= 1 and rho in (0,1/2]. Define C\_path\textasciicircum{}0 := A, C\_pol\textasciicircum{}0 := max\{B,P\}, and kappa\_pol\textasciicircum{}0 := min\{k,rho\}. Then C\_path\textasciicircum{}0,C\_pol\textasciicircum{}0 >= 1 and kappa\_pol\textasciicircum{}0 lies in (0,1/2], and all three constants are universal; this choice is made before the receiving tuple W is quantified. \steptag{validated} \\[4pt]
\step{1.2} For the witnesses fixed in 1.1, fix an arbitrary receiving Stage-1 polar witness tuple W and arbitrary algebra, delta,U\_0,U\_1,q satisfying the root hypotheses, and put g:=q+epsilon\_r*q+q\textasciicircum{}2 and r:=epsilon\_r*delta+delta\textasciicircum{}2. Then A<=C\_path, B<=C\_pol, P<=C\_pol, and kappa\_pol<=k,rho. Since epsilon\_r,delta,q are nonnegative, the receiving hypotheses imply B*(epsilon\_r+delta)<=C\_pol*(epsilon\_r+delta)<=kappa\_pol<=k, P*(epsilon\_r+delta)<=C\_pol*(epsilon\_r+delta)<=kappa\_pol<=rho, A*q<=C\_path*q<=1/4, and A*g<=C\_path*g<delta-C\_pol*r<=delta-B*r. Thus the scalar smallness and path guards required by the two cited upstream results hold. \steptag{validated} \\[4pt]
\step{1.3} Apply lem-stage1-polar-path-admissibility with its witnesses A,B,k, using exactly the transferred guards in 1.2. For every t in [0,1], L\_\{Z\_t\} is invertible and Z\_t=(1-t)U\_0+tU\_1 lies in calUbar\_\{A*g\}. Because A*g<=C\_path*g, the defining residual bound for calUbar\_\{A*g\} implies the weaker bound for calUbar\_\{C\_path*g\}, with the same right inverse; hence Z\_t lies in calUbar\_\{C\_path*(q+epsilon\_r*q+q\textasciicircum{}2)\} for every t. \steptag{validated} \\[4pt]
\step{1.4} Apply lem-stage1-polar-retraction-transport to the receiving tuple W. Its coefficient and margin hypotheses hold because P <= C\_pol and 0 < kappa\_pol <= rho by 1.1-1.2, while its algebra/delta smallness hypothesis is exactly the non-strict receiving guard C\_pol*(epsilon\_r+delta) <= kappa\_pol from the root hypotheses (also recorded in 1.2); no strict inequality is needed. For the very map displayed in the root, Pi\_delta:calU x B\_delta\textasciicircum{}\{calH\}(J)->S\_delta, Pi\_delta(U,H)=U bold-dot H, the transport theorem therefore gives a C\textasciicircum{}1 diffeomorphism onto S\_delta, with unique inverse (u\_delta,h\_delta), the identities X=u\_delta(X) bold-dot h\_delta(X), u\_delta(U)=U and h\_delta(U)=J, and the inclusion calU\_\{delta-C\_pol*r\} subseteq S\_delta. \steptag{validated} \\[4pt]
\step{1.5} For each t, 1.3 gives Z\_t in calUbar\_\{C\_path*g\}, so Z\_t has a right inverse and ||Z\_t\textasciicircum{}dagger bold-dot Z\_t-J||<=2*C\_path*g. The strict receiving guard in 1.2 gives C\_path*g<delta-C\_pol*r, hence ||Z\_t\textasciicircum{}dagger bold-dot Z\_t-J||<2*(delta-C\_pol*r). By def-approximate-unitary-space, Z\_t lies in calU\_\{delta-C\_pol*r\}; the inclusion supplied in 1.4 therefore puts Z\_t in S\_delta. Consequently u\_delta(Z\_t) in the contract is the well-typed unique first component of the inverse of that same displayed Pi\_delta. \steptag{validated} \\[4pt]
\step{1.6} The inverse (u\_delta,h\_delta) in 1.4 is C\textasciicircum{}1 and therefore continuous on S\_delta. The affine map (t,U\_0,U\_1) mapsto (1-t)U\_0+tU\_1 is jointly continuous, and 1.5 places its value in S\_delta throughout the admissible parameter domain. Therefore H(t,U\_0,U\_1):=u\_delta((1-t)U\_0+tU\_1) is jointly continuous in (t,U\_0,U\_1). \steptag{validated} \\[4pt]
\step{1.7} At t=0 and t=1 the affine path gives Z\_0=U\_0 and Z\_1=U\_1. Since U\_0,U\_1 lie in calU and 1.4 gives u\_delta(U)=U for every U in calU, H(0,U\_0,U\_1)=U\_0 and H(1,U\_0,U\_1)=U\_1; thus H joins the two endpoints. \steptag{validated} \\[4pt]
\step{1.8} Fix c in U(1) and t in [0,1], and write (u,h):=(u\_delta(Z\_t),h\_delta(Z\_t)). By 1.4-1.5, u lies in calU, h lies in B\_delta\textasciicircum{}\{calH\}(J), and Pi\_delta(u,h)=Z\_t. Conjugate-linearity of dagger, bilinearity of bold-dot, and |c|=1 show directly from def-approximate-unitary-space that cu lies in calU (its unitary residual is unchanged and a right inverse rescales by conjugate(c)); also Pi\_delta(cu,h)=c*Pi\_delta(u,h)=cZ\_t. The affine identity (1-t)cU\_0+t cU\_1=cZ\_t holds, and the same calculation puts cU\_0,cU\_1 in calU and preserves all norm guards. Since Pi\_delta is a diffeomorphism and hence has a unique inverse, its inverse at cZ\_t is (cu,h), so u\_delta(cZ\_t)=c*u\_delta(Z\_t). Therefore H(t,cU\_0,cU\_1)=c*H(t,U\_0,U\_1). \steptag{validated} \\[4pt]
\step{1.9} The constants selected in 1.1 have the required universal ranges. Because the receiving W, algebra, delta, U\_0,U\_1 and q fixed in 1.2 were arbitrary under the root hypotheses, 1.3 proves both binder-free path conclusions, 1.4-1.5 identify and type the exact displayed polar inverse along the whole path, and 1.6-1.8 prove joint continuity, the endpoint identities, and U(1)-equivariance. Universally generalizing these conclusions and existentially generalizing the constants proves the root contract. \steptag{validated} \\[4pt]
\step{1.10} Boundary-case repair for the transfers used in 1.2-1.3: with g:=q+epsilon\_r*q+q\textasciicircum{}2 and r:=epsilon\_r*delta+delta\textasciicircum{}2, nonnegativity gives g,r>=0. From A<=C\_path and B<=C\_pol one has A*g<=C\_path*g and B*r<=C\_pol*r, hence the root strict guard yields the valid mixed chain A*g<=C\_path*g<delta-C\_pol*r<=delta-B*r; in particular A*g<delta-B*r, which is exactly the strict path guard required by lem-stage1-polar-path-admissibility. Likewise kappa\_pol<=min\{k,rho\} gives only kappa\_pol<=k and kappa\_pol<=rho, and these non-strict bounds suffice: B*(epsilon\_r+delta)<=C\_pol*(epsilon\_r+delta)<=kappa\_pol<=k for path admissibility and P*(epsilon\_r+delta)<=C\_pol*(epsilon\_r+delta)<=kappa\_pol<=rho for polar-retraction transport. Thus neither C\_path=A, nor kappa\_pol=k or rho, nor g=0 causes a gap; no false strict intermediate inequality is used. \steptag{validated} \\[4pt]
\end{proofbox}

\end{document}
